\begin{quote}
\textbf{Authenticated dependencies for the preceding proof.}

\textbf{D1 (GJR Proposition~5.8).}
Let $X$ be a good strictly $K$-analytic space and let
$\Lambda\subset\mathbf R$ be $\mathbf Z$ or a divisible additive subgroup.
A function $h$ is $\Lambda$-PL if and only if every point has an analytic
neighbourhood $U$ on which
\[
                         h=\sum_{i=1}^m\lambda_i h_i,
 \qquad \lambda_i\in\Lambda,quad h_i\in\operatorname{PL}_{\mathbf Z}(U).
\]
The source proves this by choosing a $\mathbf Q$-linearly independent finite
coefficient set and gluing the integral PL pieces after shrinking.

\textbf{D2 (GJR Lemma~6.18 and the model system of Sections~6.2--6.6).}
For a $\Lambda$-linear combination of PL germs, residues are linear:
\[
                     L_{\sum_i\lambda_i h_i}(x)
                       =\sum_i\lambda_iL_{h_i}(x).
\]
Every finite family of residue line-bundle classes is represented by line
bundles on one common model $V$ of $\operatorname{red}(X,x)$, because the
model system is filtered.  Such a model is integral and has dimension
\[
             \dim V=\operatorname{trdeg}_{\widetilde K}
                         \widetilde{\mathscr H(x)}.
\]
Integer linear combinations of the representing line bundles correspond to
the same integer linear combinations of the integral PL germs.

\textbf{D3 (GJR Proposition~7.1 and Remark~7.3(4)).}
A $\Lambda$-PL germ $f$ is semipositive at $x$ if and only if its residue
class $L_f(x)$ is nef, and it is harmonic at $x$ if and only if this class is
numerically trivial.  Equivalently, after shrinking, it is represented by a
nef, respectively numerically trivial, formal model of the trivial bundle.
For any fixed PL function, the loci on which it is semipositive and on which
it is harmonic are open.

\textbf{D4 (the only algebra and curve geometry used).}
On an integral scheme of dimension at most one that is not a proper curve,
there is no integral proper closed curve, so the definition of nefness and
numerical triviality by degrees on integral proper curves is vacuous.  On a
proper integral curve, a real line-bundle class is nef if and only if its
degree is nonnegative and numerically trivial if and only if its degree is
zero.  Finally, if $d\in\mathbf Z^m$ is nonzero and
$\delta=\gcd(d_1,\ldots,d_m)$, then
$\ell=d/\delta:\mathbf Z^m\to\mathbf Z$ is primitive.  B\'ezout gives
$u$ with $\ell(u)=1$, and the exact sequence
\[
 0\longrightarrow\ker\ell\longrightarrow\mathbf Z^m
   \stackrel{\ell}{\longrightarrow}\mathbf Z\longrightarrow0
\]
splits, giving $\mathbf Z^m=\mathbf Zu\oplus\ker\ell$.  Tensoring this
integral direct sum with the additive group $\Lambda$ gives the coefficient
decomposition used in the proof; in particular
$t=\ell(\lambda)=(d\cdot\lambda)/\delta$ belongs to $\Lambda$ without any
assumption on the value group.

The local arXiv v3 source is
\texttt{Gubler\_Jell\_Rabinoff\_2111.05741.pdf}, SHA-256
\texttt{572d51c9f8907dbb814684f7d136612243c0e0b70cddbe6759bb686f0bcfbd42}.
The quoted results occur on printed pages 33, 43, and 44--46.  Immediately
before Proposition~7.5 the authors write that they do not know whether the
same integral-decomposition characterization holds for semipositive
functions.  The target theorem proves only the residual-dimension-at-most-one
positive-cone formulation; it does not assume or claim a higher-dimensional
polyhedral description of the nef cone.
\end{quote}
